% Статья о репозитории go-lgbm-serving (формат article).
% Прозу пишем отдельным заходом через skill writer; типографику - через ru-latex.
% Сборка: latexmk -pdf main.tex (движок pdflatex + biber, см. latexmkrc).
\documentclass[11pt,a4paper]{article}

\usepackage[T2A]{fontenc}
\usepackage[utf8]{inputenc}
\usepackage[russian]{babel}
\usepackage{csquotes}        % \enquote{} -> кавычки-ёлочки (см. skill writer)
\usepackage{amsmath}
\usepackage{graphicx}
\usepackage{booktabs}
\usepackage{siunitx}

% Библиография. style=numeric - переносимая заглушка; при желании заменить на
% biblatex-gost (ГОСТ Р 7.0.5) в фазе написания.
\usepackage[backend=biber,style=numeric,sorting=none]{biblatex}
\usepackage{hyperref}
\usepackage{cleveref}

\graphicspath{{figures/}}    % фигуры генерируются ноутбуками в paper/figures/
\addbibresource{assets/refs.bib}

\title{Подача древесной модели сессионного фрода из Go:\\
       битоточный паритет и нативные коды причин}
\author{Oleg Y. Danilkiff}
\date{}

\begin{document}
\maketitle

\begin{abstract}
% TODO(writer): аннотация. Источник: README "Что доказано и что нет" + DESIGN "Цель".
% Один абзац: вопрос (может ли Go подавать LightGBM-модель RBA, воспроизводить
% оценки и выдавать нативные коды причин при нужной задержке/конкуренции), приём
% (линковка той же lib_lightgbm), результат (битоточный паритет, SHAP вне горячего
% пути), граница (доказан путь подачи, не эффективность детекции).

\end{abstract}

\section{Введение}
\label{sec:intro}

% TODO(writer). Источник: docs/DESIGN.md "Цель"; README шапка.
% Центральный вопрос и четыре проверяемых подутверждения (подать / воспроизвести
% оценки / коды причин на каждое решение / задержка и конкуренция). Связующее
% ограничение - подутверждение 3 (точная атрибуция = численный паритет).

\section{Данные и признаки}
\label{sec:data}

% TODO(writer). Источник: docs/DESIGN.md "Данные и признаки"; README "Данные и
% признаки"; комментарии в training/train.py.
% Датасет RBA \cite{wiefling2022}, синтетичность, цель Is Attack IP (~9.9%),
% поведенческие RBA-признаки без утечки метки, выброс глобально-частотных
% признаков страны/ASN (98.5% gain - шорткат).

\section{Метод}
\label{sec:method}

% TODO(writer). Источник: docs/DESIGN.md "Метод"; README "Архитектура",
% "Численный паритет"; комментарии в serving/lgbm/lgbm.go.
% Приём: линковать ту же lib_lightgbm из venv Python (директива #cgo). Один
% предиктор, два вызывающих. Harness проверяет каждое подутверждение числом:
% смены решения и устойчивость топ-K кодов причин, не среднее отклонение.

\section{Рассмотренные альтернативы}
\label{sec:alternatives}

% TODO(writer). Источник: docs/DESIGN.md "Рассмотренные альтернативы".
% Двухступенчатая воронка. Фильтр 1 (семейство модели): нейросеть+post-hoc vs
% GLM vs бустинг+TreeSHAP \cite{lundberg2018,grinsztajn2022}. Фильтр 2 (путь
% подачи): LightGBM/XGBoost C API дают нативный SHAP, CatBoost-applier и ONNX
% \cite{onnx1176} - нет; чистый Go проигрывает дважды. Конкуренция: пул хэндлов,
% num_threads=1, минимальная версия LightGBM >= 3.2.0 \cite{lightgbm3771}.

\section{Результат}
\label{sec:results}

% TODO(writer). Источник: docs/DESIGN.md "Результат"; README таблицы паритета,
% производительности, качества.
% Битоточный паритет на одной сборке (maxD=0, 0 смен решения), измеренный
% кросс-платформенный остаток (SHAP <= 7.8e-16). Стоимость SHAP ~58x скоринга ->
% коды причин вне горячего пути. Качество вторично (ROC-AUC ~0.70).
% Фигуры: paper/figures/ (генерируются ноутбуками - ROC/PR, beeswarm, latency).

\section{Конвейер decline\,$\to$\,explain}
\label{sec:pipeline}

% TODO(writer). Источник: README "Конвейер decline->explain"; DESIGN "Что
% построено"; код serving/pipeline/ + serving/cmd/scorer.
% Горячий путь скорит и решает; для отклонений событие выкладывается вне горячего
% пути, нативный SHAP считает асинхронный воркер. Эндпойнты /score, /explain/{id},
% /metrics; мягкое завершение сливает очередь. TestHotPathIsolation: p99 /score
% остаётся далеко ниже одного SHAP при насыщенной очереди.

\section{Что покрыто и что вне охвата}
\label{sec:scope}

% TODO(writer). Источник: docs/DESIGN.md "Что объяснимость решений покрывает и
% что нет".
% Сделано: объяснение точное и доказуемо совпадает с поданным решением,
% устойчивость топ-K в CI. Вне охвата: валидированная таксономия adverse-action
% кодов, развёрнутый сетевой сервис, нагрузочный тест против SLA, эффективность
% детекции (модель честная, но скромная - по дизайну).

\section{Заключение}
\label{sec:conclusion}

% TODO(writer). Источник: docs/DESIGN.md ответ на центральный вопрос.
% Да: Go подаёт модель, воспроизводит оценки точно (битоточно на одной сборке, с
% измеренным остатком между архитектурами) и выдаёт точные нативные коды причин
% из того же предиктора, с устойчивостью топ-K в CI. "Можно ли" - доказано;
% "развёрнуто в продакшен" - намеренно вне охвата.


\printbibliography

\end{document}
